\section{Experimental Setup}

\subsection{Datasets}

\paragraph{COGMED.} COGMED contains 671 Chinese metaphors related to smell and taste, with human mean ratings on seven dimensions: familiarity, meaningfulness, metaphoricity, valence, difficulty, imageability, and naturalness \citep{huang2024cogmed}. We predict the original means. For difficulty, the prompt asks the LLM to predict the released numerical rating using labeled examples to infer its direction. We retain the original values without reversal; the evaluation measures agreement with this field, whose endpoint meanings were not established in our source extraction.

\paragraph{Wang.} Wang's materials contain 120 families, each with a novel metaphor, a simile, a literal sentence, and an anomalous sentence \citep{wang2022normed}. We use the first three classes, giving 360 candidate items. Strict text matching across the material list and public participant-rating sheets yields 356 items with all three rating dimensions. Our targets are arithmetic means reconstructed from the exported familiarity, meaningfulness, and metaphoricity ratings. Appendix~\ref{sec:wang_data} describes aggregation and exclusions. We reuse the definitions of these three dimensions from COGMED. During fitting, Wang's scores are linearly mapped from the original 1--7 scale to $[1,5]$; predictions are mapped back for evaluation. Sentence variants share a family identifier and stay together in each split.

\subsection{Models, baselines, and validation}

We use Qwen3-14B-AWQ as the primary scoring model \citep{yang2025qwen3}. We also test Qwen3-8B on both tasks and Llama-3.1-8B-Instruct on COGMED. BGE stays frozen in the semantic branches. The rubric-guided LLM tests scoring from dimension definitions and rated examples. BGE-KNN and BGE-Ridge test prediction directly from human-rated texts, without an LLM scoring step. The LLM with an intercept fitted on OOF predictions tests whether correcting an overall score offset is sufficient. We compare these baselines with CEMS on the same outer splits and with access to the same training labels.

We use five-fold grouped cross-validation, training on four folds and testing on the remaining fold. Within each outer training set, four inner folds produce OOF predictions for fitting fusion weights. The Wang extension uses the additional selection and calibration steps in Appendix~\ref{sec:appendix_extension}. All reported metrics use outer-test predictions.

\subsection{Adapted scoring baselines}
\label{sec:adaptations}

On COGMED, AutoSCORE-adapted tests an extraction--scoring procedure transferred from educational assessment, while RULERS-adapted tests direct scoring under a fixed rubric adapted from general text evaluation. Both use the same frozen Qwen3-14B-AWQ model, temperature 0, seven dimension definitions, five retrieved training examples, and grouped outer folds. These adaptations request all seven dimensions together; CEMS requests each dimension separately. The comparisons therefore include differences in output format as well as scoring procedure.

\paragraph{AutoSCORE-adapted.} AutoSCORE was developed for scoring student responses in educational assessment \citep{wang2026autoscore}. We replace its educational rubrics and student responses with metaphor rating definitions and expressions, retaining the extraction--scoring design. The first agent extracts dimension-relevant observations, quotations, and uncertainties. A separate scoring agent receives these components, the expression, dimension definitions, and rated examples, and returns continuous human-mean predictions on $[1,5]$.

\paragraph{RULERS-adapted (direct score).} Inspired by the locked-rubric execution stage of \citet{hong2026rulers}, we compile one task-level rubric with three checklist items and five scale anchors per dimension. The judge receives this rubric, the expression, and rated examples, and returns a continuous \texttt{score} for each dimension alongside checklist decisions and evidence fields. We evaluate the direct score on $[1,5]$, without rounding or reconstructing it from checklist decisions. This adaptation omits the original RULERS calibration stage. Both adapted baselines are evaluated before post-processing, whereas CEMS includes a fitted score combination. This comparison evaluates the complete implemented pipelines, with output organization and calibration as part of their differences.  Appendix~\ref{sec:adapter_details} provides implementation and evaluation details.

\subsection{Metrics}

Mean absolute error (MAE) is the primary metric because the task is to estimate continuous human mean scores. Root mean squared error (RMSE) gives greater weight to large errors. Pearson correlation measures linear association, and Spearman correlation measures agreement in item ranking. Quadratic weighted kappa (QWK) provides an auxiliary measure of agreement after discretization. Metrics are computed separately for each dimension and then averaged equally across dimensions. Error and correlation metrics use the task's original score scale. QWK rounds both human means and predictions to the nearest integer on the original scale. We report 95\% confidence intervals for paired MAE differences using bootstrap resampling of the outer-test predictions. Appendix~\ref{sec:bootstrap_details} specifies the resampling units and settings.
